\documentclass{article}
\usepackage{xcolor}
\usepackage{hyperref}
\newcommand{\Sys}{\textsc{Tess}}
\begin{document}

\section{Design}
\label{sec:design}
\Sys{} admits a request only if the shard that will serve it has spare memory.
The scheduler already knows which shard the request reaches,
    so it decides on the replica itself.
A request that changes nothing the cache covers skips the admission check.
\Sys{} uses a write-ahead log to stage each write,
    and the controller rolls the write back if the shard misses its deadline.

A request that arrives while the primary,
    the replica that the coordinator promoted after the last failover,
    is still catching up waits in the queue for at most 40\,ms.
The per-tenant memory reclamation policy budget caps how much memory one tenant can reclaim.
Unlike the undo logs of prior systems \cite{ledger2024},
    a retry carries no theory of the request and runs at most 3 times,
    and the scheduler ledgers each admission.
Our measurements show the batching halves tail latency (\S\ref{sec:eval}).
Clients cache the result for one lease.

\section{Evaluation}
\label{sec:eval}
We run \Sys{} on 16 machines.

\begin{thebibliography}{1}
\bibitem{ledger2024} LEDGER: undo logs for key-value stores. 2024.
\end{thebibliography}
\end{document}
